\section*{Bab \ref{ch:g10:universal-gravitation} --- Gravitasi Semesta dan Berat}

\begin{solution}{exo:g10:universal-gravitation:1}
\emph{1.} $F = \num{6.67e-11} \times \dfrac{60 \times 60}{0.80^2}
= \dfrac{\num{2.40e-7}}{0.64} \approx \qty{3.8e-7}{N}$.

\emph{2.} \omterm{def:g10:universal-gravitation:weight}{Berat} nyamuk itu
$P = \num{2.5e-6} \times 9.81 \approx \qty{2.5e-5}{N}$ --- sekitar $65$
kali \emph{lebih besar} daripada tarikan timbal balik kedua penari itu.
\omterm{def:g10:universal-gravitation:force}{Gravitasi} antara dua orang lebih lemah daripada \omterm{def:g10:universal-gravitation:weight}{berat} seekor nyamuk.
\end{solution}

\begin{solution}{exo:g10:universal-gravitation:2}
\emph{1.} $P = mg = 55 \times 9.81 \approx \qty{5.4e2}{N}$.

\emph{2.} Massanya tidak berubah: \qty{55}{kg} (banyaknya materi tetap
sama). \omterm{def:g10:universal-gravitation:weight}{Beratnya} menjadi
$P = 55 \times 1.62 \approx \qty{89}{N}$ --- enam kali lebih kecil.
\end{solution}

\begin{solution}{exo:g10:universal-gravitation:3}
\emph{1.} $g_{\text{Mars}} = \dfrac{G M}{R^2}
= \dfrac{\num{6.67e-11} \times \num{6.42e23}}{(\num{3.39e6})^2}
= \dfrac{\num{4.28e13}}{\num{1.15e13}} \approx \qty{3.73}{N/kg}$.

\emph{2.} Di Mars: $P = 900 \times 3.73 \approx \qty{3.4e3}{N}$. Di
Bumi: $P = 900 \times 9.81 \approx \qty{8.8e3}{N}$ --- penjelajah itu
sekitar $2.6$ kali lebih ringan di Mars, dengan massa yang sama.
\end{solution}

\begin{solution}{exo:g10:universal-gravitation:4}
\[
F = \num{6.67e-11} \times
\frac{\num{5.97e24} \times \num{7.35e22}}{(\num{3.84e8})^2}
= \frac{\num{2.93e37}}{\num{1.47e17}}
\approx \qty{2.0e20}{N}.
\]
Bulan mengerjakan gaya yang \emph{sama} besarnya pada Bumi, yaitu
\qty{2.0e20}{N}, dengan arah berlawanan: \omterm{def:g10:universal-gravitation:force}{gravitasi} bersifat timbal
balik.
\end{solution}

\begin{solution}{exo:g10:universal-gravitation:5}
\emph{1.} Jaraknya digandakan: $d^2$ dikalikan $4$, sehingga
$F = F_0/4$.

\emph{2.} Jaraknya dijadikan setengah: $F = 4F_0$.

\emph{3.} Kedua massanya digandakan: $m_1 m_2$ dikalikan $4$, sehingga
$F = 4F_0$.

\emph{4.} Pembilangnya $\times 3$, penyebutnya $\times 9$:
$F = F_0/3$.
\end{solution}

\begin{solution}{exo:g10:universal-gravitation:6}
\emph{1.} $F = \num{6.67e-11} \times
\dfrac{(\num{5.0e8})^2}{100^2}
= \num{6.67e-11} \times \dfrac{\num{2.5e17}}{\num{1.0e4}}
\approx \qty{1.7e3}{N}$.

\emph{2.} $m = F/g = 1667/9.81 \approx \qty{1.7e2}{kg}$: kedua kapal
itu saling menarik dengan \omterm{def:g10:universal-gravitation:weight}{berat} sebuah sepeda motor besar.

\emph{3.} $a = F/m = \num{1.7e3}/\num{5.0e8} \approx
\qty{3.3e-6}{m/s^2}$. Sekalipun tanpa penghalang, kedua tanker itu
memerlukan sekitar lima hari untuk mencapai kecepatan orang berjalan;
pada kenyataannya hambatan air dan tambatannya menenggelamkan tarikan
itu sama sekali.
\end{solution}

\begin{solution}{exo:g10:universal-gravitation:7}
\emph{1.} Neraca lengan menunjukkan \qty{3.0}{kg} di kedua dunia:
kantong dan massa bakunya sama-sama mengalami pembagian \omterm{def:g10:universal-gravitation:weight}{berat} dengan
faktor yang sama, sehingga perbandingannya tidak berubah. Neraca pegas
membaca \omterm{def:g10:universal-gravitation:weight}{beratnya}: $3.0 \times 9.81 \approx \qty{29}{N}$ di Bumi,
$3.0 \times 1.62 \approx \qty{4.9}{N}$ di Bulan.

\emph{2.} Neraca lengan: ia mengukur massa, dan massa itulah yang
dibeli pembelinya. Neraca pegas yang ditera dalam ``\omterm{def:g10:orders-of-magnitude:unit}{kilogram}'' di Bumi
akan membaca $4.9/9.81 \approx \qty{0.50}{kg}$ untuk kantong yang sama
di Bulan --- pemilik tokonya akan memberikan enam kantong dengan harga
satu.
\end{solution}

\begin{solution}{exo:g10:universal-gravitation:8}
\emph{1.} Jarak ke pusatnya:
$d = R + h = \num{6.37e6} + \num{4.0e5} = \qty{6.77e6}{m}$. Lalu
$g = \dfrac{\num{6.67e-11} \times \num{5.97e24}}{(\num{6.77e6})^2}
= \dfrac{\num{3.98e14}}{\num{4.58e13}} \approx \qty{8.69}{N/kg}$.

\emph{2.} $8.69/9.81 \approx 89\%$ dari nilainya di permukaan.

\emph{3.} Di ketinggian itu \omterm{def:g10:universal-gravitation:force}{gravitasinya} nyaris sekuat di permukaan:
antariksawan melayang karena stasiunnya dan segala isinya \emph{jatuh
bersama-sama} mengitari Bumi (jatuh bebas), bukan karena \omterm{def:g10:universal-gravitation:force}{gravitasinya}
lenyap.
\end{solution}

\begin{solution}{exo:g10:universal-gravitation:9}
Dari $g = GM/R^2$,
\[
M = \frac{g R^2}{G}
= \frac{8.87 \times (\num{6.05e6})^2}{\num{6.67e-11}}
= \frac{\num{3.25e14}}{\num{6.67e-11}}
\approx \qty{4.87e24}{kg},
\]
yaitu sekitar $0.8$ kali massa Bumi --- Venus nyaris kembaran Bumi
dalam ukuran dan massa.
\end{solution}

\begin{solution}{exo:g10:universal-gravitation:10}
\emph{1.} Matahari:
\[ F = \dfrac{\num{6.67e-11} \times \num{1.99e30}}{(\num{1.496e11})^2}
\approx \qty{5.9e-3}{N}. \]
Bulan:
\[ F = \dfrac{\num{6.67e-11} \times \num{7.35e22}}{(\num{3.84e8})^2}
\approx \qty{3.3e-5}{N}. \]

\emph{2.} Matahari menarik sekitar
$\num{5.9e-3}/\num{3.3e-5} \approx 180$ kali lebih kuat.

\emph{3.} Diameter Bumi merupakan \emph{bagian} yang jauh lebih besar
dari jarak Bumi--Bulan ($\approx 3\%$) daripada dari jarak
Bumi--Matahari ($\approx 0.01\%$), sehingga tarikan Bulan berubah jauh
lebih banyak dari sisi yang dekat ke sisi yang jauh pada bola Bumi.
Pasang surut hidup dari selisih itu --- dan di situlah Bulan menang.
\end{solution}

\begin{solution}{exo:g10:universal-gravitation:11}
\emph{1.} $g' = \dfrac{G (2M)}{(2R)^2} = \dfrac{2}{4}\,\dfrac{GM}{R^2}
= \dfrac{g}{2} \approx \qty{4.9}{N/kg}$: jari-jari yang digandakan
mengalahkan massa yang digandakan.

\emph{2.} Kita memerlukan $\dfrac{GM}{R'^2} = 2\,\dfrac{GM}{R^2}$,
sehingga $R'^2 = R^2/2$ dan $R' = R/\sqrt{2} \approx \qty{4.5e6}{m}$
--- sekitar \qty{4500}{km}.
\end{solution}

\begin{solution}{exo:g10:universal-gravitation:12}
\emph{1.} Tarikan per \omterm{def:g10:orders-of-magnitude:unit}{kilogramnya} sama:
$\dfrac{G M_{\text{Bumi}}}{x^2} = \dfrac{G M_{\text{Bulan}}}{(D-x)^2}$.
Dengan perkalian silang dan pembagian oleh $G$:
$\left(\dfrac{x}{D-x}\right)^2 = \dfrac{M_{\text{Bumi}}}{M_{\text{Bulan}}}$.

\emph{2.} $\dfrac{x}{D-x}
= \sqrt{\dfrac{\num{5.97e24}}{\num{7.35e22}}} = \sqrt{81.2} \approx
9.0$, sehingga $x = 9.0\,(D - x)$, yaitu $x = \dfrac{9.0}{10.0}\,D
\approx 0.90 \times \num{3.84e8} \approx \qty{3.5e8}{m}$. Titik
setimbangnya berada pada $90\%$ perjalanan menuju Bulan: tarikan Bumi
menguasai hampir seluruh perjalanan itu.
\end{solution}

\begin{solution}{exo:g10:universal-gravitation:13}
\emph{1.} Dijabarkan: $R^2 + 2Rh + h^2 = R^2 + x^2$; karena $h$ sangat
kecil dibandingkan $R$, buanglah $h^2$: $2Rh \approx x^2$, sehingga
$h \approx \dfrac{x^2}{2R}
= \dfrac{(\num{8.0e3})^2}{2 \times \num{6.37e6}}
= \dfrac{\num{6.4e7}}{\num{1.27e7}} \approx \qty{5.0}{m}$.

\emph{2.} Dalam satu sekon peluru itu jatuh \qty{4.9}{m} --- hampir
persis sebesar \qty{5.0}{m} yang ditinggalkan tanah sepanjang
\qty{8.0}{km}. Jadi pada $v \approx \qty{8.0}{km/s}$ jatuhnya tidak
pernah menyusul tanahnya: peluru itu mengorbit.

\emph{3.} $T = \dfrac{2\pi R}{v}
= \dfrac{2\pi \times \num{6.37e6}}{\num{8.0e3}}
\approx \qty{5.0e3}{s} \approx \qty{83}{min}$ --- dekat dengan sekitar
$\qty{90}{min}$ milik \omterm{def:g10:universal-gravitation:satellite}{satelit} \omterm{def:g10:universal-gravitation:satellite}{orbit} rendah yang sebenarnya, yang
terbang sedikit lebih tinggi.
\end{solution}

\begin{solution}{exo:g10:universal-gravitation:14}
\emph{1.} $g = \dfrac{\num{6.67e-11} \times \num{2.8e30}}
{(\num{1.2e4})^2} = \dfrac{\num{1.87e20}}{\num{1.44e8}}
\approx \qty{1.3e12}{N/kg}.$

\emph{2.} $P = 70 \times \num{1.3e12} \approx \qty{9.1e13}{N}$,
berbanding $\qty{6.9e2}{N}$ di Bumi: sekitar $\num{1.3e11}$ kali lebih
besar. Tidak ada bangunan yang tersusun dari atom sanggup bertahan
berdiri di sana.

\emph{3.} $P = \num{1.0e-6} \times \num{1.3e12} \approx
\qty{1.3e6}{N}$. Di Bumi itu adalah \omterm{def:g10:universal-gravitation:weight}{berat}
$m = \num{1.3e6}/9.81 \approx \qty{1.3e5}{kg}$ --- sebutir pasir yang
seberat lokomotif seratus ton.
\end{solution}

\begin{solution}{exo:g10:universal-gravitation:15}
\emph{1.} Jaraknya $60$ kali lebih besar, sehingga tarikan per
\omterm{def:g10:orders-of-magnitude:unit}{kilogramnya} $60^2 = 3600$ kali lebih lemah.

\emph{2.} Jatuhnya dalam satu sekon berskala dengan cara yang sama:
$4.9/3600 \approx \qty{1.4e-3}{m}$ --- sekitar \qty{1.4}{mm}.

\emph{3.} $T = 27.3 \times \qty{86400}{s} = \qty{2.36e6}{s}$, sehingga
$v = \dfrac{2\pi \times \num{3.84e8}}{\num{2.36e6}} \approx
\qty{1.02e3}{m/s}$, dan $x = \qty{1.02e3}{m}$ dalam satu sekon.
Jatuhnya adalah $h \approx \dfrac{x^2}{2d}
= \dfrac{(\num{1.02e3})^2}{2 \times \num{3.84e8}}
\approx \qty{1.4e-3}{m}$.

\emph{4.} Kedua bilangan itu sesuai: tiap sekon Bulan jatuh ke arah
Bumi persis sebanyak yang diramalkan \omterm{def:g10:universal-gravitation:force}{gravitasi} berkebalikan kuadrat
yang ditera pada sebuah apel yang jatuh. Newton menyimpulkan bahwa gaya
yang menahan Bulan \emph{adalah} gaya yang menjatuhkan apel ---
\omterm{def:g10:universal-gravitation:force}{gravitasi} bersifat semesta.
\end{solution}

\begin{solution}{pb:g10:universal-gravitation:1}
\textbf{1.} $F = \num{6.67e-11} \times
\dfrac{\num{7.3e6} \times 55}{100^2} \approx \qty{2.7e-6}{N}$.

\textbf{2.} Butir pasir itu \omterm{def:g10:universal-gravitation:weight}{beratnya}
$\num{1.0e-6} \times 9.81 \approx \qty{9.8e-6}{N}$: seluruh Menara
Eiffel menarik si pelancong sekitar $3.7$ kali \emph{lebih lemah}
daripada \omterm{def:g10:universal-gravitation:weight}{berat} sebutir pasir.

\textbf{3.} (a) Jaraknya $\times 10$: gayanya dibagi $100$.
(b) Kedua massanya $\times 1000$: gayanya dikalikan $10^6$.
(c) Massanya $\times 1000$ dan jaraknya $\times 1000$:
$10^6 / 10^6$ --- gayanya tidak berubah.

\textbf{4.} $F = \dfrac{\num{6.67e-11} \times \num{5.97e24} \times
1.00}{(\num{6.37e6})^2} \approx \qty{9.81}{N}$: \omterm{def:g10:universal-gravitation:weight}{berat} satu \omterm{def:g10:orders-of-magnitude:unit}{kilogram}
--- inilah tepatnya kuat medan $g = \qty{9.81}{N/kg}$.

\textbf{5.} \omterm{def:g10:universal-gravitation:force}{Gravitasi} lemah per \omterm{def:g10:orders-of-magnitude:unit}{kilogram} tetapi hanya pernah
\emph{menarik}, sehingga tarikan seluruh $\num{5.97e24}$ \omterm{def:g10:orders-of-magnitude:unit}{kilogram} Bumi
bertambah pada arah yang sama --- dan massa berukuran astronomi
mengalahkan massa seorang tetangga dengan dua puluh dua pangkat
sepuluh.

\textbf{6.} $F = \num{6.67e-11} \times \dfrac{158 \times 0.73}{0.23^2}
= \num{6.67e-11} \times \dfrac{115.3}{0.0529}
\approx \qty{1.5e-7}{N}$.

\textbf{7.} Bola kecil itu \omterm{def:g10:universal-gravitation:weight}{beratnya}
$0.73 \times 9.81 \approx \qty{7.2}{N}$ --- sekitar $\num{5e7}$ kali
tarikan yang hendak dideteksi. Tidak ada neraca berpiring yang sanggup
memisahkan satu bagian dari lima puluh juta; kawat puntir yang halus,
yang terpuntir secara kasatmata di bawah momen gaya berukuran
nanonewton, sanggup melakukannya.

\textbf{8.} $G = \dfrac{F d^2}{m_1 m_2}
= \dfrac{\num{1.47e-7} \times 0.0529}{115.3}
\approx \qty{6.74e-11}{N.m^2/kg^2}$ --- meleset sekitar $1\%$ dari
nilai modern \num{6.67e-11}, dari sebuah gudang kayu pada tahun 1798.
\textbf{9.} Dari $g = G M / R^2$:
\[
M = \frac{g R^2}{G}
= \frac{9.81 \times (\num{6.37e6})^2}{\num{6.67e-11}}
\approx \qty{5.97e24}{kg}.
\]
Enam juta miliar miliar \omterm{def:g10:orders-of-magnitude:unit}{kilogram}: Bumi, tertimbang.

\textbf{10.} $V = \frac43 \pi R^3 = \frac43 \pi
(\num{6.37e6})^3 \approx \qty{1.08e21}{m^3}$, sehingga
$\rho = \num{5.97e24}/\num{1.08e21} \approx \qty{5.5e3}{kg/m^3}$ ---
dua kali massa jenis batuan permukaan. Bagian dalamnya tentu
menyembunyikan sesuatu yang jauh lebih rapat daripada granit: Bumi
memiliki inti yang berat, yaitu inti besi.

\textbf{11.} $g_{\text{Bulan}} = \dfrac{\num{6.67e-11} \times
\num{7.35e22}}{(\num{1.74e6})^2} = \dfrac{\num{4.90e12}}
{\num{3.03e12}} \approx \qty{1.62}{N/kg}$.

\textbf{12.} Tingginya $\propto 1/g$:
$h = 0.50 \times \dfrac{9.81}{1.62} \approx \qty{3.0}{m}$ --- lompatan
di tempat yang melewati ring bola basket, dalam gerak lambat.

\textbf{13.} $g_{\text{Merkurius}} = \dfrac{\num{6.67e-11} \times
\num{3.30e23}}{(\num{2.44e6})^2} = \dfrac{\num{2.20e13}}
{\num{5.95e12}} \approx \qty{3.70}{N/kg}$ --- hampir persis sama dengan
\qty{3.73}{N/kg} milik Mars. Massa Merkurius lebih kecil, tetapi
jari-jarinya yang lebih kecil ($R^2$ yang lebih kecil di penyebut)
mengimbanginya: Merkurius adalah dunia yang jauh lebih rapat.

\textbf{14.} Kuat medan setengahnya menuntut
$(R + h)^2 = 2 R^2$, sehingga $R + h = R\sqrt{2}$ dan
$h = (\sqrt{2} - 1) R \approx 0.414 \times \num{6.37e6}
\approx \qty{2.6e6}{m}$ --- sekitar \qty{2600}{km} di atas permukaan.

\textbf{15.} $\dfrac{G M_{\text{Bumi}}}{d^2} = 1.62$ memberikan
$d = \sqrt{\dfrac{\num{6.67e-11} \times \num{5.97e24}}{1.62}}
= \sqrt{\num{2.46e14}} \approx \qty{1.57e7}{m} \approx 2.5$ jari-jari
Bumi. Hanya dari satu setengah jari-jari di atas tanah, tarikan Bumi
tidak lebih kuat daripada tarikan di permukaan Bulan.

\textbf{16.} $v = \dfrac{2\pi d}{T}
= \dfrac{2\pi \times \num{1.496e11}}{\num{3.156e7}}
\approx \qty{2.98e4}{m/s}$ --- tiga puluh kilometer tiap sekon.

\textbf{17.} $\dfrac{v^2}{d} = \dfrac{(\num{2.98e4})^2}{\num{1.496e11}}
\approx \qty{5.9e-3}{N/kg}$: Matahari menarik tiap \omterm{def:g10:orders-of-magnitude:unit}{kilogram} Bumi dengan
sekitar enam milinewton.

\textbf{18.} $\dfrac{G M_{\text{Matahari}}}{d^2} = \dfrac{v^2}{d}$
memberikan
\[
M_{\text{Matahari}} = \frac{v^2 d}{G}
= \frac{\num{8.87e8} \times \num{1.496e11}}{\num{6.67e-11}}
\approx \qty{1.99e30}{kg}.
\]

\textbf{19.} $\dfrac{\num{1.99e30}}{\num{5.97e24}} \approx
\num{3.3e5}$: Matahari setara sepertiga juta Bumi.

\textbf{20.} Karena $G$ yang sama mengatur setiap pasang massa,
mengukurnya sekali saja antara dua bola timbal di sebuah gudang
memungkinkan kita menimbang mula-mula planet di bawah kaki kita, lalu
bintang yang kita orbiti --- itulah harga kata \emph{semesta}.
\end{solution}
